\documentclass[11pt]{article}

\usepackage[margin=0.82in]{geometry}
\usepackage{amsmath,amssymb,bm}
\usepackage{microtype}
\usepackage[hidelinks]{hyperref}
\usepackage{enumitem}
\usepackage{xcolor}
\usepackage{booktabs}

\newcommand{\e}{\epsilon}
\newcommand{\ii}{\mathrm{i}}
\newcommand{\bG}{\mathbf{G}}
\newcommand{\bQ}{\mathbf{Q}}
\newcommand{\bK}{\mathbf{K}}
\newcommand{\bk}{\mathbf{k}}
\newcommand{\br}{\mathbf{r}}
\newcommand{\bu}{\mathbf{u}}
\newcommand{\cO}{\mathcal{O}}

\title{\vspace{-1.2cm}Charge-Density-Wave Diffraction from a Commensurate Supercell}
\author{}
\date{}

\begin{document}
\maketitle
\vspace{-0.9cm}

Consider a simple cubic lattice with a static displacement modulation
\begin{equation}
    \bu(\br)=\bm{\e}\sin(\bQ\cdot\br),
    \qquad
    \bm{\e}\parallel[100],
    \qquad
    \bQ\parallel[100].
\end{equation}
For a commensurate modulation of wavelength $Ma$,
\begin{equation}
    \bQ=\frac{2\pi}{Ma}\hat{\mathbf{x}},
\end{equation}
so one CDW supercell contains $M$ primitive cells along $x$.  Within one such supercell,
\begin{equation}
    \br_j^0=ja\,\hat{\mathbf{x}},\qquad j=0,\ldots,M-1,
\end{equation}
and the displaced coordinates are
\begin{equation}
    \br_j=\br_j^0+\bm{\e}\sin(\bQ\cdot\br_j^0).
\end{equation}
For identical atoms of form factor $f$, the relevant supercell structure factor is
\begin{equation}
    F(\bk)=f\sum_{j=0}^{M-1}e^{\ii\bk\cdot\br_j}.
    \label{eq:Fexact}
\end{equation}
The unmodulated $y$ and $z$ directions simply impose the usual primitive-lattice reciprocal conditions; the nontrivial selection rule is along $x$.

Assuming $\e\ll a$, expand the phase factor,
\begin{align}
 e^{\ii\bk\cdot\br_j}
 &=e^{\ii\bk\cdot\br_j^0}
   e^{\ii\bk\cdot\bm{\e}\sin(\bQ\cdot\br_j^0)}\\
 &\simeq e^{\ii\bk\cdot\br_j^0}
 \left[1+\ii(\bk\cdot\bm{\e})\sin(\bQ\cdot\br_j^0)\right].
\end{align}
Thus
\begin{equation}
    F(\bk)=F_0(\bk)+F_1(\bk)+\cO(\e^2),
\end{equation}
with
\begin{align}
F_0(\bk)
&=f\sum_{j=0}^{M-1}e^{\ii k_xja},\\
F_1(\bk)
&=\ii f(\bk\cdot\bm{\e})
  \sum_{j=0}^{M-1}e^{\ii k_xja}
  \sin\!\left(\frac{2\pi j}{M}\right).
\end{align}

Index the supercell reciprocal points along $x$ by
\begin{equation}
    k_x=\frac{2\pi}{Ma}n_x.
\end{equation}
Then
\begin{equation}
F_0=f\sum_{j=0}^{M-1}e^{\ii2\pi n_xj/M}
=
\begin{cases}
Mf,&n_x=0\pmod M,\\
0,&\text{otherwise}.
\end{cases}
\end{equation}
Hence the ordinary cubic Bragg peaks occur at $n_x=mM$; in the lowest nonzero zone the parent peak is $n_x=M$.

For the first-order term, using
\begin{equation}
\ii\sin\theta=\frac{1}{2}\left(e^{\ii\theta}-e^{-\ii\theta}\right),
\end{equation}
one obtains
\begin{align}
F_1
=\frac{f(\bk\cdot\bm{\e})}{2}
\left[
\sum_{j=0}^{M-1}e^{\ii2\pi(n_x+1)j/M}
-
\sum_{j=0}^{M-1}e^{\ii2\pi(n_x-1)j/M}
\right].
\end{align}
Each roots-of-unity sum vanishes unless its exponent is an integer multiple of $2\pi$. Therefore
\begin{equation}
    F_1\neq0
    \quad\Longleftrightarrow\quad
    n_x=\pm1\pmod M.
\end{equation}
Around the parent peak $n_x=M$, the only first-order CDW satellites are therefore
\begin{equation}
    \boxed{n_x=M-1\quad\text{and}\quad n_x=M+1},
\end{equation}
or equivalently
\begin{equation}
    \boxed{\bk=\bG-\bQ\quad\text{and}\quad\bk=\bG+\bQ}.
\end{equation}

Their structure factors are
\begin{equation}
    \boxed{
    F(\bG-\bQ)
    \simeq
    +\frac{Mf}{2}(\bk\cdot\bm{\e})
    },
\end{equation}
\begin{equation}
    \boxed{
    F(\bG+\bQ)
    \simeq
    -\frac{Mf}{2}(\bk\cdot\bm{\e})
    },
\end{equation}
up to the global Fourier-sign convention. Since $\bm{\e}\parallel\hat{\mathbf{x}}$,
\begin{equation}
    \bk\cdot\bm{\e}=k_x\e.
\end{equation}
Consequently the leading satellite intensity is
\begin{equation}
    \boxed{I_{\pm Q}\propto |F_1|^2\propto(\bk\cdot\bm{\e})^2\propto\e^2.}
\end{equation}

So the $M$-site supercell sum cancels all new reciprocal points to first order except the two roots-of-unity conditions $n_x=M\pm1$, giving the CDW satellites $\bG\pm\bQ$ with amplitude $\pm(Mf/2)(\bk\cdot\bm{\e})$ and intensity proportional to $\e^2$.

\newpage
% ============================================================
\textbf{Additional note: higher-order satellites (often seen experimentally)}\\
% ============================================================

In fact, a perfectly sinusoidal CDW already produces $2Q,3Q,\ldots$ peaks. Higher-order satellites do \emph{not} automatically imply an anharmonic CDW. Even for the perfectly sinusoidal displacement
\begin{equation}
    u_x(x)=\e_1\sin(Qx),
\end{equation}
the scattering phase $e^{\ii K_xu_x}$ is nonlinear in the displacement. Higher order diffraction peaks can occur with larger displacement amplitudes despite perfect periodicity. Using the Jacobi--Anger expansion,
\begin{equation}
    e^{\ii z\sin\theta}
    =\sum_{m=-\infty}^{\infty}J_m(z)e^{\ii m\theta},
\end{equation}
Eq.~\eqref{eq:Fexact} gives satellites at
\begin{equation}
    \boxed{\bK=\bG+m\bQ,\qquad m=0,\pm1,\pm2,\ldots}
\end{equation}
with, for a single sinusoidal displacement and $Q\ll G$,
\begin{equation}
    F_m\simeq Mf\,J_m(x),
    \qquad
    I_m\simeq M^2|f|^2J_m^2(x),
    \qquad
    x\equiv|\bG\cdot\bm{\e}_1|.
    \label{eq:Bessel}
\end{equation}
(The exact argument is the scattering vector of the individual satellite, $|\bK_m\cdot\bm{\e}_1|$; replacing it by $|\bG\cdot\bm{\e}_1|$ is appropriate when $Q\ll G$.)

For $x\ll1$,
\begin{equation}
    J_m(x)\simeq\frac{1}{m!}\left(\frac{x}{2}\right)^m,
\end{equation}
so that
\begin{equation}
    \boxed{
    I_m\propto
    \frac{(x/2)^{2m}}{(m!)^2}.
    }
    \label{eq:factorial}
\end{equation}
In particular,
\begin{align}
 I_1&\propto \frac{x^2}{4},\\
 I_2&\propto \frac{x^4}{64},\\
 I_3&\propto \frac{x^6}{2304},
\end{align}
and therefore
\begin{equation}
    \boxed{\frac{I_2}{I_1}\simeq\frac{x^2}{16}},
    \qquad
    \boxed{\frac{I_3}{I_2}\simeq\frac{x^2}{36}},
\end{equation}
with the general adjacent-order ratio
\begin{equation}
    \boxed{
    \frac{I_{m+1}}{I_m}
    \simeq
    \frac{x^2}{4(m+1)^2}.
    }
    \label{eq:adjacentratio}
\end{equation}
Thus the rapid, approximately factorial decay $Q\to2Q\to3Q\cdots$ is the baseline expectation for a weak, purely sinusoidal periodic lattice distortion (PLD).

One may then use the harmonic decay to measure the displacement amplitude. The intensity hierachy above provides an internal measurement of $\e_1$.  Instead of inferring the displacement only from the weak first-order satellite, one can fit several orders simultaneously to Eq.~\eqref{eq:Bessel}.  For example, in the weak limit,
\begin{equation}
    x\simeq4\sqrt{\frac{I_2}{I_1}}
    \simeq6\sqrt{\frac{I_3}{I_2}},
\end{equation}
and therefore
\begin{equation}
    \e_1\simeq\frac{x}{|\bG\cdot\hat{\bm{\e}}|}.
\end{equation}
Agreement of the values inferred from $I_1$, $I_2/I_1$, $I_3/I_2$, and the depletion of the parent Bragg intensity through $J_0^2(x)$ is a strong self-consistency test of the simple sinusoidal-displacement model.

In practice, the integrated intensities should first be corrected for the ordinary diffraction factors (atomic form factor, polarization/Lorentz factors as appropriate, absorption, Debye--Waller factor, illuminated volume, detector geometry, and domain population).  Comparing harmonics around the \emph{same} parent $\bG$ minimizes many of these systematic factors.

We can then also discuss concrete physical information encoded in $Q:2Q:3Q:\cdots$.

\subsection*{Nonlinear coupling and secondary order parameters}
Let $\psi_Q$ be the primary CDW order parameter and let $\eta_{2Q}$ denote a possible second-harmonic structural/electronic order parameter.  Translational symmetry can permit a Landau coupling of the form
\begin{equation}
    \mathcal{F}_{\mathrm{int}}
    =\lambda_2\eta_{2Q}^*\psi_Q^2+\mathrm{c.c.}
\end{equation}
(and similarly $\eta_{3Q}^*\psi_Q^3$, or mixed terms involving $\psi_Q\eta_{2Q}$).  If the higher harmonic is a secondary, ``slaved'' order parameter, minimizing the free energy gives schematically
\begin{equation}
    \eta_{2Q}\propto\psi_Q^2,
    \qquad
    \eta_{3Q}\propto\psi_Q^3.
\end{equation}
After appropriate structure-factor corrections this implies
\begin{equation}
    I_{2Q}\propto I_Q^2,
    \qquad
    I_{3Q}\propto I_Q^3,
\end{equation}
or slopes of approximately 2 and 3 in log--log plots of $I_{mQ}$ versus $I_Q$.

However, the pure Bessel/kinematic mechanism has the \emph{same leading powers} in the small-displacement limit.  The discriminator is therefore the \emph{coefficient and evolution}, not the exponent alone.  One should compare the measured $I_{mQ}$ with the parameter-free higher-order intensity predicted after fixing $\e_1$ from the fundamental peak.  Additional harmonic order is indicated when the higher-order intensity cannot be reproduced by that Bessel baseline.

A useful experimental precedent is that in blue bronze $\mathrm{K}_{0.3}\mathrm{MoO}_3$, the approximate ratio $I_{m=2}:I_{m=1}:I_{\mathrm{main}}\sim10^{-4}:10^{-2}:1$ remains nearly temperature independent below $T_{\mathrm{CDW}}$, consistent with second-order scattering from an essentially sinusoidal PLD.  In $\mathrm{NbSe}_3$, by contrast, the second-order satellite grows anomalously strongly upon cooling relative to the first-order satellite, which has been interpreted as development of a second-harmonic/secondary order parameter.\cite{vanSmaalen2005}

\subsection*{Real-space waveform and ``squaring up''}
The diffraction structural factor is essentially a fourier decomposition of the real space waveform. Once the purely kinematic Bessel contribution is removed, the residual hierarchy $|\e_1|,|\e_2|,|\e_3|,\ldots$ reconstructs how sinusoidal or non-sinusoidal the CDW is.

A weak, smooth sinusoid has essentially only $\e_1$.  Strong commensurate lock-in or a soliton/discommensuration texture can sharpen the real-space modulation and produce a much slower harmonic decay.  In the limiting intuition of a square-like waveform, Fourier amplitudes decay only algebraically (e.g. $\sim1/m$ for the allowed harmonics), so the corresponding weak-displacement diffraction intensities decay roughly as $1/m^2$ rather than factorially.  Thus a conspicuously slow $Q\to2Q\to3Q$ decay is a direct signature that the real-space modulation has acquired sharp features or strong anharmonicity.

\subsection*{Symmetry information from even versus odd harmonics}
A displacement waveform satisfying the half-period antisymmetry
\begin{equation}
    u(x+\pi/Q)=-u(x)
\end{equation}
has no \emph{intrinsic} even Fourier components $\e_2,\e_4,\ldots$.  Consequently, a genuine even-harmonic structural component can diagnose loss of that waveform symmetry or coupling to an additional order parameter.

There is an important caveat: even-order \emph{diffraction satellites} are still generated by the Bessel expansion of a purely sinusoidal $u(x)$.  Therefore symmetry should be inferred from the extracted Fourier components of the displacement after the kinematic contribution is modeled, not only from the observation of a $2Q$ peak.

\subsection*{Distinguishing a harmonic from an independent density wave}
A peak near $2Q$ need not be a harmonic at all; it can belong to a second independent ordering vector.  Useful tests are:
\begin{itemize}[leftmargin=1.6em]
    \item whether its wavevector remains locked to exactly $2Q(T)$ as temperature, pressure, field, or doping changes;
    \item whether its onset temperature coincides with the primary CDW or appears at a separate transition;
    \item whether its linewidth/correlation length tracks that of the $Q$ peak;
    \item whether its intensity follows the Bessel/secondary-order scaling expected from $I_Q$.
\end{itemize}
Failure of these tests points toward an independent mode rather than a simple harmonic.  Diffraction studies of CDW materials contain examples of both possibilities.\cite{vanSmaalen2005}

\textbf{Higher harmonics can, in this sense, also as a probe of fluctuations near a phase transition. }Near a continuous CDW transition, the static ordered displacement normally becomes \emph{small}; $\e\sim a$ is not the generic critical limit.  The important physics is instead the spatial and temporal fluctuation of the complex order parameter,
\begin{equation}
    \psi(\br)=|\psi(\br)|e^{\ii\phi(\br)},
\end{equation}
and the displacement may be written schematically as
\begin{equation}
    \bu(\br)=\mathrm{Re}\left[\bm{\psi}(\br)e^{\ii\bQ\cdot\br}\right].
\end{equation}
Finite correlation length produces diffuse/broadened satellite scattering with
\begin{equation}
    \Delta q\sim\xi^{-1}.
\end{equation}
Diffuse CDW scattering above $T_{\mathrm{CDW}}$ and sharpening into long-range superlattice peaks on cooling are standard signatures of this regime.\cite{vanSmaalen2005}

Higher harmonics can be especially sensitive to phase disorder.  If a local $m$th harmonic carries phase $m\phi(\br)$, its correlation function contains
\begin{equation}
    \left\langle e^{\ii m[\phi(\br)-\phi(0)]}\right\rangle.
\end{equation}
For approximately Gaussian phase fluctuations,
\begin{equation}
    \left\langle e^{\ii m\Delta\phi}\right\rangle
    =\exp\left[-\frac{m^2}{2}\left\langle(\Delta\phi)^2\right\rangle\right].
\end{equation}
Thus $2Q$ and $3Q$ coherence can be suppressed more rapidly than the fundamental by phason/phase fluctuations.  A higher-order peak that broadens or loses integrated coherent weight much faster than the $Q$ peak can therefore reveal phase fluctuations even when the local distortion amplitude remains sizable.  The exact relation between harmonic linewidths and $m$ is model dependent, but the $m\phi$ phase sensitivity is general.

\begin{thebibliography}{9}

\bibitem{Overhauser1971}
A.~W. Overhauser,
``Observability of Charge-Density Waves by Neutron Diffraction,''
\emph{Phys. Rev. B} \textbf{3}, 3173 (1971).
\href{https://doi.org/10.1103/PhysRevB.3.3173}{doi:10.1103/PhysRevB.3.3173}.

\bibitem{Overhauser1986}
A.~W. Overhauser,
``Satellite-intensity patterns from the charge-density wave in potassium,''
\emph{Phys. Rev. B} \textbf{34}, 7362 (1986).
\href{https://doi.org/10.1103/PhysRevB.34.7362}{doi:10.1103/PhysRevB.34.7362}.

\bibitem{vanSmaalen2005}
S.~van Smaalen,
``The Peierls transition in low-dimensional electronic crystals,''
\emph{Acta Crystallographica Section A} \textbf{61}, 51--61 (2005).
\href{https://doi.org/10.1107/S0108767304028414}{doi:10.1107/S0108767304028414}.

\bibitem{Chang2016}
J.~Chang \emph{et al.},
``The microscopic structure of charge density waves in underdoped YBa$_2$Cu$_3$O$_{6.54}$ revealed by X-ray diffraction,''
\emph{Nature Communications} \textbf{7}, 11494 (2016).
\href{https://doi.org/10.1038/ncomms11494}{doi:10.1038/ncomms11494}.

\end{thebibliography}

\end{document}
